% ECONOMICS_PROBLEM: {"id": "D63-130", "title": "Fair allocation across changing profiles", "jel_code": "D63", "source_id": "OP-0278"}
% !TeX program = xelatex
\documentclass[11pt,a4paper]{article}
\usepackage[margin=23mm]{geometry}
\usepackage{fontspec}
\setmainfont{Latin Modern Roman}
\usepackage{amsmath,amssymb,mathtools}
\usepackage{xcolor,enumitem,booktabs,longtable,array,xurl,hyperref}
\definecolor{muted}{HTML}{53636C}
\hypersetup{colorlinks=true,urlcolor=blue,linkcolor=blue}
\setlength{\parindent}{0pt}
\setlength{\parskip}{5pt}
\newcommand{\R}{\mathbb R}
\newcommand{\E}{\mathbb E}
\newcommand{\N}{\mathbb N}
\newcommand{\ind}{\mathbf 1}
\newcommand{\Var}{\operatorname{Var}}
\newcommand{\Cov}{\operatorname{Cov}}
\newcommand{\supp}{\operatorname{supp}}
\DeclareMathOperator*{\argmin}{arg\,min}
\DeclareMathOperator*{\argmax}{arg\,max}
\newcommand{\entrymeta}[3]{{\sffamily\small\color{muted}\textbf{\texttt{#1}}\quad|\quad #2\par #3\par}}
\newcommand{\Problemref}[1]{\texttt{#1}}
\newcommand{\JELref}[2]{\texttt{#2} (JEL \texttt{#1})}
\begin{document}
\section*{Fair allocation across changing profiles}
\textbf{Problem ID:} \texttt{D63-130}\quad \textbf{JEL:} \texttt{D63}\par
% BEGIN ECONOMICS STATEMENT
\label{problem:OP-0278}
{\sffamily\small\color{muted}Original subject group: Economic theory, equilibrium and social choice\par}
\label{definitions:problem:30007560}
\textbf{Audit classification:} Published research agenda. Dagstuhl §4.5 explicitly proposes inter-profile axioms. The four-axiom classification is editorial; a fixed finite domain is decidable by finite search.
\subsection*{Definitions and precise research target}
\textbf{Model and research target.} Fix a finite universe of agents and goods and integer \(V\geq1\). Each profile consists of agent set \(N\), good set \(G\), and additive valuations \(v_i(g)\in\{0,\ldots,V\}\). A deterministic allocation rule \(F\) selects a complete allocation for every profile in a specified domain \(\mathcal D\). It is envy-free up to one good if, for each ordered pair \(i,j\), either \(F_j=\varnothing\) or some \(g\in F_j\) satisfies \(v_i(F_i)\geq v_i(F_j\setminus\{g\})\). It is Pareto efficient if no other complete allocation weakly raises every agent's value and strictly raises at least one. Resource monotonicity requires that adding a good without changing old valuations weakly increases every incumbent's value. Population monotonicity requires that adding an agent without changing incumbent valuations weakly decreases every incumbent's value. Determine which subsets of these four requirements are jointly attainable on \(\mathcal D\), and characterize the maximal domains or explicit weakenings that permit compatibility. Quantifiers must range over every admissible profile and each permitted profile transition, rather than only selected examples. Supply constructive rules or impossibility proofs, with the relevant valuation and population restrictions made explicit.

\emph{Definitions and maintained assumptions (editorial unless a source definition is named).} Let \(\mathcal D\) be a schema generating profiles at all finite population and resource sizes, with integer values at most the input parameter \(V\). A rule family chooses \(F(P)\) for every generated profile. A resource transition \(P\to P'\) preserves agents and all old item values and adds exactly one new item; a population transition preserves goods and incumbent values and adds one new agent. Define each monotonicity inequality using the incumbent's unchanged valuation. EF1 compares every ordered pair and deletes one item from the envied bundle; PO excludes any feasible allocation improving all values weakly and one strictly. Define domains by recognizable valuation restrictions, such as identical or binary additive valuations, and close them under the transitions under study. For one fully fixed finite universe and \(V\), enumeration of rule outputs and their cross-profile constraints decides compatibility. The nontrivial task is a structural theorem for a domain schema of unbounded size, an explicit impossibility witness, or an efficiently implementable rule family.
\subsection*{Variants and assumption changes}
\begin{enumerate}
\item \textbf{Exact axioms (editorial specialization).} Classify EF1, PO and resource monotonicity on identical, binary and unrestricted additive domains across all sizes; report already known constructions separately from any residual gap.
\item \textbf{Randomized transitions (editorial).} Replace value inequalities by expected-value inequalities under a randomized rule, keeping EF1 either ex post or in an explicitly defined weaker expectation sense. Determine which deterministic impossibilities survive; those two fairness interpretations are distinct.
\end{enumerate}
\subsection*{References and source audit}
\begin{enumerate}
\item Dagstuhl §4.5. Supports the stated research area; exact displayed specification is editorial. \url{https://drops.dagstuhl.de/storage/04dagstuhl-reports/volume14/issue09/24401/html/DagRep.14.9.145/DagRep.14.9.145.html}
\end{enumerate}
\begin{enumerate}\raggedright
\item\hypertarget{definitions:source:30007560:1}{} Dominik Peters. 2024. \emph{Inter-profile Axioms in the Allocation of Indivisible Goods, in Fair Division: Algorithms, Solution Concepts, and Applications}. Seminar report §4.5.
\url{https://drops.dagstuhl.de/storage/04dagstuhl-reports/volume14/issue09/24401/html/DagRep.14.9.145/DagRep.14.9.145.html}.
DOI: \url{https://doi.org/10.4230/DagRep.14.9.145}.
\end{enumerate}

\subsection*{Common conventions: Source mathematical conventions}

These conventions apply to each entry unless explicitly replaced.

\textbf{Models and identification.} A primitive model \(m\) specifies all probability laws, feasible choices, information, equilibrium and measurement rules used by its target. The admissible class \(\mathcal M\) is fixed independently of outcomes. If \(P_O(m)\) is its observable probability law and \(\tau(m)\) the target, its identified set at observable law \(P\) is
\[
 \mathcal I_\tau(P)=\{\tau(m):m\in\mathcal M,\ P_O(m)=P\}.
\]
Point identification means that this set is a singleton. An empty set signals incompatibility of data and assumptions. Coordinate bounds need not describe the joint set. Bounds are sharp only if every retained target is attainable or arbitrarily approachable by compatible models. Finite bounds require explicit restrictions on unobserved quantities; unrestricted confounding, hidden wealth, extrapolation or arbitrary equilibrium selection can yield uninformative sets.

\textbf{Causal interventions.} A potential outcome \(Y_i(p)\) is the outcome under a complete policy regime \(p\), including timing, financing, information and market spillovers. The target population and units are fixed before treatment. With interference, use the complete assignment vector or a justified exposure mapping; individual no-interference notation alone cannot identify an economy-wide intervention. A difference of mean potential outcomes does not require the cross-world joint distribution. Conditioning on post-treatment migration, survival, enrollment or employment changes the estimand and needs separate justification.

\textbf{Statistical statements.} An estimator, confidence set, forecast or test specifies a sampling law, model class and loss/coverage criterion. Uniform \((1-\alpha)\)-coverage means \(\inf_{m\in\mathcal M}\Pr_m\{\tau(m)\in C_n\}\geq1-\alpha\), or an explicitly asymptotic version. The whole parameter space trivially has coverage, so informativeness requires a stated width, risk or power objective. A forecasting improvement is relative to a named benchmark, information set and evaluation population.

\textbf{Equilibrium and optimization.} An equilibrium comprises feasible plans, optimality, beliefs where needed, budget constraints and all market/resource clearing. The primitives or a stated benchmark define each correspondence \(\mathcal E(p;m)\). Nonemptiness is assumed or proved, never inferred just from compact policy sets. A selected-equilibrium target uses a declared selection rule; a robust target quantifies over all permitted equilibria. Use suprema or infima unless attainment is established. An \(\varepsilon\)-optimal policy is feasible and within \(\varepsilon\) of the finite optimum value. Report an empty feasible set separately.

\textbf{Welfare and accounting.} Normative utility, interpersonal normalization, social weights, numeraire, discounting and the use of public receipts are primitives. Behavioral utility need not coincide with the welfare criterion. Household welfare includes ownership income and rents once; adding profits again to compensating variation generally double-counts them. A monetary equivalent is defined by an explicit transfer and price convention, with monotonicity and range conditions. Average effects, mechanism contributions, transfers and welfare losses are different targets. Infinite sums require convergence and dynamic feasible plans require terminal or no-Ponzi conditions.

\textbf{Source versions.} A working-paper date, revision date and journal publication year are distinguished. A DOI for a journal version is not automatically the DOI of an earlier manuscript. Locators refer to the stated version and printed pages where specified. A metadata-only check does not verify a claim about a full-text conclusion. Every entry's source subsection narrows the provenance claim accordingly.
% END ECONOMICS STATEMENT
\end{document}
